% ----------------------------------------------------------------
% AMS-LaTeX Paper ************************************************
% **** -----------------------------------------------------------
%\documentclass{amsart}
%\usepackage{txfonts}
%\documentclass[12pt,oneside]{article}
\documentclass{amsart}
\usepackage{graphicx}
\usepackage{enumitem}
% ----------------------------------------------------------------
\vfuzz2pt % Don't report over-full v-boxes if over-edge is small
\hfuzz2pt % Don't report over-full h-boxes if over-edge is small
% THEOREMS -------------------------------------------------------
\newtheorem{thm}{Theorem}[section]
\newtheorem{cor}[thm]{Corollary}
\newtheorem{lem}[thm]{Lemma}
\newtheorem{prop}[thm]{Proposition}
\theoremstyle{definition}
\newtheorem{defn}[thm]{Definition}
\theoremstyle{Exercise}
\newtheorem{ex}[thm]{Exercise}
\theoremstyle{remark}
\newtheorem{rem}[thm]{Remark}
\theoremstyle{rule}
\newtheorem{rul}[thm]{Rule}

\numberwithin{equation}{section}
% MATH -----------------------------------------------------------
\newcommand{\norm}[1]{\left\Vert#1\right\Vert}
\newcommand{\abs}[1]{\left\vert#1\right\vert}
\newcommand{\set}[1]{\left\{#1\right\}}
\newcommand{\Real}{\mathbb R}
\newcommand{\Z}{\mathbb Z}
\newcommand{\To}{\longrightarrow}
\newcommand{\BX}{\bB(X)}
\newcommand{\A}{\mathcal{A}}
% ----------------------------------------------------------------

% define some simple, commonly-used commands
\newcommand{\eps}{\varepsilon}
\newcommand{\dsum}{\displaystyle\sum}
\newcommand{\dint}{\displaystyle\int}

\newcommand{\pdr}[2]{\dfrac{\partial{#1}}{\partial{#2}}}
\newcommand{\pdrr}[2]{\dfrac{\partial^2{#1}}{\partial{#2}^2}}
\newcommand{\pdrt}[3]{\dfrac{\partial^2{#1}}{\partial{#2}{\partial{#3}}}}
\newcommand{\dr}[2]{\dfrac{d{#1}}{d{#2}}}
\newcommand{\aver}[1]{\langle {#1} \rangle}
\newcommand{\Baver}[1]{\Big\langle {#1} \Big\rangle}

\newcommand{\bzero}{\mathbf 0}
\newcommand{\bGamma}{\mbox{\boldmath{$\Gamma$}}}
\newcommand{\btheta}{\boldsymbol \theta}
\newcommand{\bchi}{\mbox{\boldmath{$\chi$}}}
\newcommand{\bnu}{\boldsymbol \nu}
\newcommand{\bmu}{\boldsymbol \mu}
\newcommand{\brho}{\mbox{\boldmath{$\rho$}}}
\newcommand{\bxi}{\boldsymbol \xi}
\newcommand{\bnabla}{\boldsymbol \nabla}
\newcommand{\bOm}{\boldsymbol \Omega}
\newcommand{\blambda}{\boldsymbol \lambda}
\newcommand{\bsigma}{\boldsymbol \sigma}

\newcommand{\bbR}{\mathbb R}
\newcommand{\bbC}{\mathbb C}
\newcommand{\bbQ}{\mathbb Q}
\newcommand{\bbN}{\mathbb N}
\newcommand{\bbZ}{\mathbb Z}

\newcommand{\ba}{\mathbf a} \newcommand{\bb}{\mathbf b}
\newcommand{\bc}{\mathbf c} \newcommand{\bd}{\mathbf d}
\newcommand{\be}{\mathbf e} \newcommand{\bff}{\mathbf f}
\newcommand{\bg}{\mathbf g} \newcommand{\bh}{\mathbf h}
\newcommand{\bi}{\mathbf i} \newcommand{\bj}{\mathbf j}
\newcommand{\bk}{\mathbf k} \newcommand{\bl}{\mathbf l}
\newcommand{\bm}{\mathbf m} \newcommand{\bn}{\mathbf n}
\newcommand{\bo}{\mathbf o} \newcommand{\bp}{\mathbf p}
\newcommand{\bq}{\mathbf q} \newcommand{\br}{\mathbf r}
\newcommand{\bs}{\mathbf s} \newcommand{\bt}{\mathbf t}
\newcommand{\bu}{\mathbf u} \newcommand{\bv}{\mathbf v}
\newcommand{\bw}{\mathbf w} \newcommand{\bx}{\mathbf x}
\newcommand{\by}{\mathbf y} \newcommand{\bz}{\mathbf z}
\newcommand{\bA}{\mathbf A} \newcommand{\bB}{\mathbf B}
\newcommand{\bC}{\mathbf C} \newcommand{\bD}{\mathbf D}
\newcommand{\bE}{\mathbf E} \newcommand{\bF}{\mathbf F}
\newcommand{\bG}{\mathbf G} \newcommand{\bH}{\mathbf H}
\newcommand{\bI}{\mathbf I} \newcommand{\bJ}{\mathbf J}
\newcommand{\bK}{\mathbf K} \newcommand{\bL}{\mathbf L}
\newcommand{\bM}{\mathbf M} \newcommand{\bN}{\mathbf N}
\newcommand{\bO}{\mathbf O} \newcommand{\bP}{\mathbf P}
\newcommand{\bQ}{\mathbf Q} \newcommand{\bR}{\mathbf R}
\newcommand{\bS}{\mathbf S} \newcommand{\bT}{\mathbf T}
\newcommand{\bU}{\mathbf U} \newcommand{\bV}{\mathbf V}
\newcommand{\bW}{\mathbf W} \newcommand{\bX}{\mathbf X}
\newcommand{\bY}{\mathbf Y} \newcommand{\bZ}{\mathbf Z}

\newcommand{\cA}{\mathcal A} \newcommand{\cB}{\mathcal B}
\newcommand{\cC}{\mathcal C} \newcommand{\cD}{\mathcal D}
\newcommand{\cE}{\mathcal E} \newcommand{\cF}{\mathcal F}
\newcommand{\cG}{\mathcal G} \newcommand{\cH}{\mathcal H}
\newcommand{\cI}{\mathcal I} \newcommand{\cJ}{\mathcal J}
\newcommand{\cK}{\mathcal K} \newcommand{\cL}{\mathcal L}
\newcommand{\cM}{\mathcal M} \newcommand{\cN}{\mathcal N}
\newcommand{\cO}{\mathcal O} \newcommand{\cP}{\mathcal P}
\newcommand{\cQ}{\mathcal Q} \newcommand{\cR}{\mathcal R}
\newcommand{\cS}{\mathcal S} \newcommand{\cT}{\mathcal T}
\newcommand{\cU}{\mathcal U} \newcommand{\cV}{\mathcal V}
\newcommand{\cW}{\mathcal W} \newcommand{\cX}{\mathcal X}
\newcommand{\cY}{\mathcal Y} \newcommand{\cZ}{\mathcal Z}


%%%%%%%%%%%%%%Start%%%%%%%%%%%%%Start%%%%%%%%%%%Start%%%%%%%%%%%%%%%Start%%%%%%%%%%%%%%%%%%%%%%%%%Start%%%%%%%%%%%%%%%%
\usepackage{fancyhdr}

\pagestyle{fancy}
\fancyhf{}
\rhead{}
\chead{\includegraphics[scale=.1]{snhu_logo.png}}
\begin{document}

\title{\sf Module One Problem Set}%

%\thm{bbjh}


\begin{center}
\includegraphics[scale=.1]{snhu_logo.png}
\end{center}

\maketitle 
This document is proprietary to Southern New Hampshire University. It and the problems within may not be posted on any non-SNHU website.
\\\\\\\\
\begin{center}
%Enter your name below this line:
Tyler Ellis
\end{center}

\begin{center}
\rule{\textwidth}{0.4pt}
\end{center}
\newpage
%--------------------------------------------------------------------------------------------------

\section*{}

\section*{}
Directions: Type your solutions into this document and be sure to show all steps for arriving at your solution. Just giving a final number may not receive full credit.
\\

\section*{Problem 1}
\begin{enumerate}[label=(\alph*)]

\item Every patient was given the medication or the placebo or both.\\

% logical translation
\[
  \forall x\,\bigl(D(x)\lor P(x)\bigr).
\]

% negation + De Morgan
\[
  \neg\forall x\,(D(x)\lor P(x))
  \;\equiv\;
  \exists x\,\neg\bigl(D(x)\lor P(x)\bigr)
  \;\equiv\;
  \exists x\,\bigl(\neg D(x)\land \neg P(x)\bigr).
\]

% English rendering of the negation
\textit{English: There is at least one patient who was given neither the medication nor the placebo.}

\item Every patient who took the placebo had migraines. (Hint: you will need to apply the conditional identity, $p \to q \equiv \neg p \lor q$.)\\

% logical translation using p→q ≡ ¬p∨q
\[
  \forall x\;\bigl(P(x)\to M(x)\bigr)
  \;\equiv\;
  \forall x\;\bigl(\neg P(x)\lor M(x)\bigr).
\]

% negation + De Morgan
\[
  \neg\forall x\,(\neg P(x)\lor M(x))
  \;\equiv\;
  \exists x\,\neg\bigl(\neg P(x)\lor M(x)\bigr)
  \;\equiv\;
  \exists x\,\bigl(P(x)\land \neg M(x)\bigr).
\]

% English rendering of the negation
\textit{English: There is at least one patient who took the placebo and did not have migraines.}

\item There is a patient who had migraines and was given the placebo.\\

% logical translation
\[
  \exists x\,\bigl(M(x)\land P(x)\bigr).
\]

% negation + De Morgan
\[
  \neg\exists x\,(M(x)\land P(x))
  \;\equiv\;
  \forall x\,\neg\bigl(M(x)\land P(x)\bigr)
  \;\equiv\;
  \forall x\,\bigl(\neg M(x)\lor \neg P(x)\bigr).
\]

% English rendering of the negation
\textit{English: Every patient either did not have migraines or was not given the placebo (or both).}

\end{enumerate}

\newpage
%--------------------------------------------------------------------------------------------------

\section*{Problem 2}
Use De Morgan's law for quantified statements and the laws of propositional logic to show the following equivalences:
\begin{enumerate}[label=(\alph*)]

\item 
\[
  \neg \forall x \,\bigl(P(x)\land \neg Q(x)\bigr)
  \;\equiv\;
  \exists x\,\neg\bigl(P(x)\land \neg Q(x)\bigr)
  \;\equiv\;
  \exists x\,\bigl(\neg P(x)\lor Q(x)\bigr).
\]

\item 
\[
  \neg \forall x \,\bigl(\neg P(x)\to Q(x)\bigr)
  \;\equiv\;
  \neg \forall x \,\bigl(P(x)\lor Q(x)\bigr)
  \;\equiv\;
  \exists x\,\neg\bigl(P(x)\lor Q(x)\bigr)
  \;\equiv\;
  \exists x\,\bigl(\neg P(x)\land \neg Q(x)\bigr).
\]

\item 
\[
  \neg \exists x \,\bigl(\neg P(x)\lor (Q(x)\land \neg R(x))\bigr)
  \;\equiv\;
  \forall x\,\neg\bigl(\neg P(x)\lor (Q(x)\land \neg R(x))\bigr)
  \;\equiv\;
  \forall x\,\bigl(P(x)\land (\neg Q(x)\lor R(x))\bigr).
\]

\end{enumerate}

\newpage
%--------------------------------------------------------------------------------------------------

\section*{Problem 3}
The domain of discourse for this problem is a group of three people who are working on a project. To make notation easier, the people are numbered $1,2,3$. The predicate $M(x,y)$ indicates whether $x$ has sent an email to $y$. The table below shows the value of the predicate.

\[
 \begin{array}{||c||c|c|c||}
\hline\hline
M & 1 & 2 & 3\\
\hline\hline
1 & T & T & T\\
\hline
2 & T & F & T\\
\hline
3 & T & T & F\\
\hline\hline
\end{array}
\]

Determine if each quantified statement is true or false. Justify your answer.

\begin{enumerate}[label=(\alph*)]

\item 
\[
  \forall x\,\forall y\;\bigl(x\neq y\to M(x,y)\bigr).
\]
\textbf{Answer:} True.  
\textbf{Justification:} For every pair $x\neq y$ among $\{1,2,3\}$, $M(x,y)=T$ in the table, so each implication holds.

\item 
\[
  \forall x\,\exists y\;\neg M(x,y).
\]
\textbf{Answer:} False.  
\textbf{Justification:} For $x=1$, $M(1,1)=T$, $M(1,2)=T$, $M(1,3)=T$, so there is no $y$ with $\neg M(1,y)$. Thus the “$\exists y$” fails for $x=1$.

\item 
\[
  \exists x\,\forall y\;M(x,y).
\]
\textbf{Answer:} True.  
\textbf{Justification:} For $x=1$, $M(1,1)=M(1,2)=M(1,3)=T$, so there exists an $x$ (namely 1) for which $M(x,y)$ holds for all $y$.

\end{enumerate}

\newpage
%--------------------------------------------------------------------------------------------------

\section*{Problem 4}
Translate each English statement into logical expressions. The domain of discourse is the set of all real numbers.

\begin{enumerate}[label=(\alph*)]

\item The reciprocal of every positive number less than one is greater than one.\\
\[
  \forall x\,\bigl((0<x\land x<1)\to \tfrac1x>1\bigr).
\]

\item There is no smallest number.\\
\[
  \neg\exists m\,\forall y\,(m\le y).
\]
% or equivalently \(\forall x\,\exists y\,(y<x)\).

\item Every number other than 0 has a multiplicative inverse.\\
\[
  \forall x\,\bigl(x\neq 0\to \exists y\,(x\cdot y=1)\bigr).
\]

\end{enumerate}

\newpage
%--------------------------------------------------------------------------------------------------

\section*{Problem 5}
The sets $A$, $B$, and $C$ are defined as follows:

\[
  A = \{\text{tall},\;\text{grande},\;\text{venti}\},\quad
  B = \{\text{foam},\;\text{no‑foam}\},\quad
  C = \{\text{non‑fat},\;\text{whole}\}.
\]

Use the definitions for $A$, $B$, and $C$ to answer the questions. Express the elements using $n$-tuple notation.

\begin{enumerate}[label=(\alph*)]
  \item Write an element from the set $A\times B\times C$.\\
  \(\;(\text{tall},\,\text{foam},\,\text{whole})\)

  \item Write an element from the set $B\times A\times C$.\\
  \(\;(\text{no‑foam},\,\text{venti},\,\text{non‑fat})\)

  \item Write the set $B\times C$ using roster notation.\\
  \[
    B\times C
    =\{\,
      (\text{foam},\,\text{non‑fat}),\,
      (\text{foam},\,\text{whole}),\,
      (\text{no‑foam},\,\text{non‑fat}),\,
      (\text{no‑foam},\,\text{whole})
    \}.
  \]
\end{enumerate}

\end{document}